激光超声波检测相比传统接触式探头具有显著优势：它能够实现非接触检测、支持高速扫描，且可应用于复杂几何形状和粗糙表面，无需耦合剂\cite{liu2022deep}。这些特性使激光超声波在航空航天、制造和过程控制环境中的无损检测中具有特别的吸引力。
然而，由于激光烧蚀过程的随机性和热弹性转换特性，激光产生的超声波信号存在信噪比(SNR)较低的问题\cite{zhang2020deep}。产生的噪声掩盖了诊断相关的声学特征——包括波达时间、幅度比和波形形态——这些是可靠缺陷检测和表征的关键\cite{chen2019wavelet}。

信号平均是激光超声波中信噪比改进最广泛采用的方法\cite{zhang2020deep}。通过对同一测量进行多次采集和叠加，随机噪声分量相互抵消，而相干信号逐渐增强，从而获得与平均次数平方根成正比的SNR增益。
尽管有效，但大量平均也带来了实际限制：采集时间随平均次数线性增加，重复激光脉冲可能导致敏感或易碎材料的表面损伤，且在高速度工业检测中变得不切实际\cite{yang2023unsupervised}。这些限制促使我们寻求去噪方法，能够在不依赖大量平均的情况下实现相当或更优的信噪比改善。

传统去噪技术——包括维纳滤波、小波阈值和块匹配协同滤波(BM3D)——可以在一定程度上抑制噪声\cite{dabov2007bm3d, chen2019wavelet}。然而，这些方法通常针对均方误差(MSE)或SNR改善等整体噪声指标进行优化，缺乏对保留声学有意义特征的显式约束\cite{yang2023unsupervised}。对于缺陷检测，诊断特征的保真度——频散模式的精确到达时间、反射回波的相对幅度以及波形包络的形状——与整体噪声抑制同等重要。
在抑制噪声的同时破坏这些特征几乎没有任何实际价值。
本工作旨在解决如何实现有效的激光超声波信号去噪，同时明确保留对缺陷检测至关重要的声学特征，从而减少对大量信号平均的依赖。